% section: パターン別解法

% 既存の「基本4パターン」「応用7パターン」を,実際の見抜き方まで
% 含めて整理した節である.

\section{パターン別解法}

\subsection*{この節の読み方}

第2章では漸化式の基本的な型を学んだ.ここでは,それらを実際の問題で
どう見抜き,どの順に変形するかをまとめる.まず目の前の式を見て,
\[
  \text{この形なら,何を置けば既知の型になるか}
\]
を判断しなければならない.

この節では,各パターンについて,次の順に説明する.

\begin{enumerate}
  \item どこを見ればその型だと分かるか.
  \item 何を目標に式変形するか.
  \item 置換や公式を使って,実際にどう解くか.
  \item その解法を,どのような一般原理として理解できるか.
\end{enumerate}

ここで扱う「型」は,問題文に書いてある名前ではない.式の中に現れる
構造を読み取り,既知の型へ翻訳する方法である.

\subsection{パターン1：定数項がある一次漸化式は,不動点を引く}

\subsubsection{この形なら何をするか}

\[
  a_{n+1}=pa_n+q
\]
のように,\,$a_n$ に定数倍を施した後で定数を足している場合を考える.
等比型
\[
  b_{n+1}=pb_n
\]
に近いが,\,$q$ が邪魔をしている.

このときは,数列からある定数 $\alpha$ を引いて,
\[
  b_n=a_n-\alpha
\]
と置く.
うまくいく $\alpha$ は,
\[
  \alpha=p\alpha+q
\]
を満たす数,つまり写像 $x\mapsto px+q$ の不動点である.

\noindent\textbf{【方針】}

\begin{enumerate}
  \item $\alpha=p\alpha+q$ を解き,不動点を求める.
  \item $b_n=a_n-\alpha$ と置く.
  \item $b_{n+1}=pb_n$ という等比型に帰着する.
\end{enumerate}

\subsubsection{例題}

\[
  a_1=1,\qquad a_{n+1}=3a_n-4
\]
の一般項を求める.

\noindent\textbf{【解法】}

不動点 $\alpha$ は
\[
  \alpha=3\alpha-4
\]
を満たすから,
\[
  \alpha=2
\]
である.
そこで
\[
  b_n=a_n-2
\]
と置く.
すると,
\[
  \begin{aligned}
    b_{n+1}
     & =a_{n+1}-2 \\
     & =3a_n-4-2  \\
     & =3(a_n-2)  \\
     & =3b_n.
  \end{aligned}
\]
よって $b_n$ は等比型であり,
\[
  b_n=b_1\,3^{n-1}=(1-2)3^{n-1}=-3^{n-1}.
\]
したがって,
\begin{mdquote}
  $a_n=2-3^{n-1}$
\end{mdquote}

\subsubsection{なぜこの置き方を思いつくのか}

定数を引くと定数項が消えるのは偶然ではない.
右辺の $pa_n+q$ が,\,$a_n=\alpha$ をそのまま返す点を中心に,
\[
  a_{n+1}-\alpha=p(a_n-\alpha)
\]
と書き直したからである.

この一手で,数列の値そのものではなく,
不動点からの距離が何倍されるかを調べる問題に変わった.
これは第2章の特性方程式型の解法であり,付録では固定点と安定性の理論へ発展する.

\subsection{パターン2：右辺に $f(n)$ が残るなら,補助数列を設計する}

\subsubsection{等比型にしたいという目標を先に決める}

\[
  a_{n+1}=pa_n+f(n)
\]
では,定数項ではなく $n$ の関数 $f(n)$ が現れる.
そこで,ある関数 $g(n)$ を加えて,
\[
  b_n=a_n+g(n)
\]
と置く.
目標は,
\[
  b_{n+1}=pb_n
\]
である.

実際に必要な条件を書き下すと,
\[
  a_{n+1}+g(n+1)=p\{a_n+g(n)\}.
\]
元の漸化式を代入すれば,
\[
  f(n)+g(n+1)=pg(n)
\]
となる.
つまり,\,$g(n)$ は
\[
  g(n+1)-pg(n)=-f(n)
\]
を満たすように設計する.

\noindent\textbf{【見抜き方】}

\begin{itemize}
  \item $f(n)$ が一次式なら,まず $g(n)$ も一次式と予想する.
  \item $f(n)$ が二次式なら,まず $g(n)$ も二次式と予想する.
  \item 予想した形を代入し,係数比較で未知係数を決める.
\end{itemize}

\subsubsection{例題：$a_{n+1}=2a_n+n$}

\[
  a_1=0,\qquad a_{n+1}=2a_n+n
\]
を考える.
右辺の余分な項は $n$ だから,\,$g(n)=An+B$ と予想する.
\[
  g(n+1)-2g(n)=-n
\]
に代入すると,
\[
  \{A(n+1)+B\}-2(An+B)
  =-An+(A-B).
\]
これが $-n$ に一致するには,
\[
  A=1,\qquad A-B=0
\]
なので,
\[
  g(n)=n+1
\]
とすればよい.

実際に
\[
  b_n=a_n+n+1
\]
と置くと,
\[
  \begin{aligned}
    b_{n+1}
     & =a_{n+1}+(n+1)+1 \\
     & =2a_n+n+n+2      \\
     & =2(a_n+n+1)      \\
     & =2b_n.
  \end{aligned}
\]
初項は $b_1=2$ だから,
\[
  b_n=2^n.
\]
したがって,
\begin{mdquote}
  $a_n=2^n-n-1$
\end{mdquote}

\subsubsection{発展：補助数列は「特解を足す」方法である}

$g(n)$ を探す方法は,単なる技巧ではない.
元の漸化式のうち,\,$f(n)$ を打ち消す特別な数列を一つ作り,
それを使って同次型の解へ移っている.

線形漸化式では,
\[
  \text{一般解}
  =
  \text{同次漸化式の解}
  +
  \text{特別な一つの解}
\]
という考え方が繰り返し現れる.
この「余分な項を消すための補助解」という見方が,
線形代数や微分方程式の特解へつながる.

\subsection{パターン3：差が $f(n)$ なら,和を取る}

\subsubsection{階差型の見抜き方}

\[
  a_{n+1}=a_n+f(n)
\]
を見たら,
\[
  a_{n+1}-a_n=f(n)
\]
と左辺を差の形にする.
すると,隣り合う式を書き並べたときに途中の項が消える.

\noindent\textbf{【方針】}

\[
  \begin{aligned}
    a_2-a_1     & =f(1),   \\
    a_3-a_2     & =f(2),   \\
                & \ \vdots \\
    a_n-a_{n-1} & =f(n-1)
  \end{aligned}
\]
を足し合わせる.
左辺は隣り合う項が打ち消し合い,
\[
  (a_2-a_1)+(a_3-a_2)+\cdots+(a_n-a_{n-1})
  =a_n-a_1
\]
となる.
したがって,
\begin{mdquote}
  $a_n=a_1+\sum_{k=1}^{n-1}f(k)$
\end{mdquote}

\subsubsection{例題：差が二次式の場合}

\[
  a_1=1,\qquad a_{n+1}=a_n+n^2
\]
なら,
\[
  a_n=1+\sum_{k=1}^{n-1}k^2
  =1+\frac{(n-1)n(2n-1)}6.
\]

ここで,和の公式をただ暗記するのではなく,
\[
  S_n-S_{n-1}=n^2
\]
となる三次式 $S_n$ を設計することもできる.
\[
  S_n=An^3+Bn^2+Cn+D
\]
と置き,差を取って係数比較する.
差分は多項式の次数を一つ下げるので,右辺が二次式なら左辺の候補は三次式になる.

\subsubsection{差分から離散積分へ}

微分に対して積分が逆操作であるように,
差分に対して和が逆操作になる.
\[
  \Delta a_n=a_{n+1}-a_n=f(n)
\]
を解くことは,
\[
  \Delta^{-1}f(n)
\]
を求めることだと考えられる.
この見方を進めると,有限差分法,離散的な微積分,
さらには母関数へ進む.

\subsection{パターン4：隣の項との比が分かるなら,積を整理する}

\subsubsection{階比型の見抜き方}

\[
  a_{n+1}=f(n)a_n
\]
では,項の差ではなく比
\[
  \frac{a_{n+1}}{a_n}=f(n)
\]
が分かっている.
そのため,最初から
\[
  a_n=a_1\prod_{k=1}^{n-1}f(k)
\]
と書くのが基本である.

ただし,積がそのままでは扱いにくい場合がある.
そこで,\,$f(n)$ を
\[
  f(n)=\frac{g(n)}{g(n+1)}
\]
と表せれば,積が望ましく消える.

\subsubsection{例題：望ましい積を作る}

\[
  a_1=1,\qquad
  a_{n+1}=\frac{n}{n+1}a_n
\]
を考える.
係数は
\[
  \frac{n}{n+1}
\]
だから,
\[
  a_n
  =\prod_{k=1}^{n-1}\frac{k}{k+1}
  =\frac1n.
\]

同じことを「型をそろえる」方法で書くと,
\[
  b_n=na_n
\]
と置けば,
\[
  b_{n+1}
  =(n+1)a_{n+1}
  =(n+1)\frac{n}{n+1}a_n
  =na_n=b_n.
\]
したがって $b_n=1$ であり,\,$a_n=1/n$ となる.

\subsubsection{発展：積を和に変える}

係数が複雑なときは,対数を取ることで
\[
  \log a_{n+1}-\log a_n=\log f(n)
\]
とできる.
積の問題が差の問題に変わる.
正値数列の階比型が,階差型やガンマ関数の漸化式へつながる理由はここにある.

\subsection{パターン5：指数が数列にかかっているなら,対数を取る}

\subsubsection{べき型の見抜き方}

\[
  a_{n+1}=p\,a_n^q
\]
のように $a_n$ がべきになっている場合,\,$a_n$ のままでは指数が繰り返し現れる.
正の数列であることを確認したうえで,
\[
  b_n=\log a_n
\]
と置く.
すると,
\[
  b_{n+1}
  =\log p+q\log a_n
  =q b_n+\log p
\]
となる.
対数によって,乗法とべきが一次漸化式へ変わった.

\subsubsection{例題}

\[
  a_1=1,\qquad a_{n+1}=4a_n^2
\]
を考える.
底を $2$ とする対数
\[
  b_n=\log_2a_n
\]
を取ると,
\[
  b_{n+1}=2+2b_n.
\]
これは特性方程式型である.
不動点は $-2$ なので,
\[
  b_n+2=2(b_{n-1}+2).
\]
初項 $b_1=0$ から,
\[
  b_n+2=2^n,
  \qquad
  b_n=2^n-2.
\]
したがって,
\begin{mdquote}
  $a_n=2^{\,2^n-2}$
\end{mdquote}

\subsubsection{底は何を選んでもよい}

対数の底は本質ではない.
底を $e$ にしても,底を $2$ にしても,一次漸化式になる.
ただし,\,$p$ が $2$ のべきなら底を $2$ にする方が係数が整数になりやすい.

解法の中で大切なのは,
\[
  a_n\text{ の指数を外したい}
\]
という目的から対数を選ぶことである.

\subsection{パターン6：逆数を取ると分数式が消える}

\subsubsection{分数型の見抜き方}

\[
  a_{n+1}=\frac{a_n}{1+ca_n}
\]
のように,分母にも $a_n$ が現れる場合,
分数を直接展開するより,逆数を取る.
\[
  \frac1{a_{n+1}}
  =\frac{1+ca_n}{a_n}
  =\frac1{a_n}+c.
\]

\noindent\textbf{【方針】}

\[
  b_n=\frac1{a_n}
\]
とおき,\,$b_n$ の階差型に帰着する.

\subsubsection{例題}

\[
  a_1=1,\qquad
  a_{n+1}=\frac{a_n}{1+2a_n}
\]
を考える.
逆数を取ると,
\[
  \frac1{a_{n+1}}
  =\frac1{a_n}+2.
\]
そこで $b_n=1/a_n$ とおくと,
\[
  b_{n+1}=b_n+2,\qquad b_1=1.
\]
よって
\[
  b_n=2n-1
\]
であり,
\begin{mdquote}
  $a_n=\frac1{2n-1}$
\end{mdquote}

\subsubsection{分数一次変換への入口}

\[
  a_{n+1}=\frac{\alpha a_n+\beta}{\gamma a_n+\delta}
\]
のような場合は,単純な逆数だけでは足りない.
しかし,固定点 $\lambda,\mu$ を見つけて
\[
  b_n=\frac{a_n-\lambda}{a_n-\mu}
\]
と置くと,\,$b_n$ が等比型になることがある.
逆数変換は,分数一次変換を単純化する変数変換の最初の例である.

\subsection{パターン7：三角関数型は,倍角公式を漸化式へ翻訳する}

\subsubsection{この形なら何を思い出すか}

\[
  a_{n+1}=2\cos\theta\,a_n-a_{n-1}
\]
の形を見たら,次の倍角・加法公式を思い出す.
\[
  \cos((n+1)\theta)
  =2\cos\theta\cos(n\theta)-\cos((n-1)\theta).
\]

したがって,
\[
  a_n=\cos(n\theta)
\]
という形を予想できる.
ここで大切なのは,公式を見て答えを暗記することではない.
漸化式の右辺が,三角関数の加法公式の右辺と同じ形をしているから,
\[
  a_n=\cos\theta_n,\qquad \theta_n=n\theta
\]
と置くのである.

\subsubsection{例題}

\[
  a_0=1,\qquad a_1=\cos\theta,\qquad
  a_{n+1}=2\cos\theta\,a_n-a_{n-1}
\]
を考え,
\[
  a_n=\cos(n\theta)
\]
を示す.

\noindent\textbf{【方針】}

\begin{enumerate}
  \item 倍角・加法公式から,候補 $\cos(n\theta)$ が同じ漸化式を満たすことを確認する.
  \item 初期値 $a_0,a_1$ が候補と一致することを確認する.
  \item 以上の二つから,帰納法で全ての $n$ について一致させる.
\end{enumerate}

\noindent\textbf{【解法】}

まず,
\[
  \cos(0\theta)=1=a_0,\qquad
  \cos(1\theta)=\cos\theta=a_1
\]
である.
次に,\,$a_{n-1}=\cos((n-1)\theta)$,\,$a_n=\cos(n\theta)$ が成り立つと仮定する.
このとき,
\[
  \begin{aligned}
    a_{n+1}
     & =2\cos\theta\,a_n-a_{n-1} \\
     & =2\cos\theta\cos(n\theta)
    -\cos((n-1)\theta)           \\
     & =\cos((n+1)\theta).
  \end{aligned}
\]
最後の等号は
\[
  2\cos\theta\cos(n\theta)
  =\cos((n+1)\theta)+\cos((n-1)\theta)
\]
を用いたものである.
よって帰納法により,
\begin{mdquote}
  $a_n=\cos(n\theta)$
\end{mdquote}
が示された.

\subsubsection{なぜ帰納法が自然なのか}

二階漸化式では,次の項を決めるために直前二項が必要である.
したがって,
\[
  n=0,1\text{ で一致する}
\]
だけでは不十分で,二項が一致しているとき次の項も一致することを示す必要がある.
三角関数の候補が漸化式を満たすことを確認し,
初期値を合わせれば,候補と数列の一致が永遠に保たれる.

これは「一般項を思いついたから帰納法で確認する」というより,
「同じ更新規則と同じ初期状態を持つ二つの列は同じである」
という一意性を使っている.

\subsection{パターン8：振動型は,\,$\cos(n\theta)$ と $\sin(n\theta)$ を組み合わせる}

\subsubsection{初期値が一般の場合}

\[
  a_{n+1}=2\cos\theta\,a_n-a_{n-1}
\]
で,
\[
  a_0=A,\qquad a_1=B
\]
とする.
この場合,単に $\cos(n\theta)$ だけでは初期値を調整できない.
そこで,
\[
  a_n=C\cos(n\theta)+D\sin(n\theta)
\]
と予想する.

この二つの列は,それぞれ同じ漸化式を満たす.
したがって,それらの線形結合も同じ漸化式を満たす.
初期値から
\[
  C=A,\qquad
  A\cos\theta+D\sin\theta=B
\]
を決める.\,$\sin\theta\neq0$ なら,
\[
  D=\frac{B-A\cos\theta}{\sin\theta}.
\]

\subsubsection{複素数による一行の説明}

オイラーの公式
\[
  e^{i\theta}=\cos\theta+i\sin\theta
\]
を用いると,
\[
  z^n,\qquad z=e^{i\theta}
\]
が特性方程式
\[
  z^2-2\cos\theta\,z+1=0
\]
の根になる.
実部と虚部を取れば,\,$\cos(n\theta)$ と $\sin(n\theta)$ が現れる.

つまり,三角関数型は,
\[
  \text{倍角公式による帰納法}
\]
と
\[
  \text{複素数の特性方程式}
\]
という二つの説明を持っている.
高校生向けには前者から始め,後者を発展として示すと流れがよい.

\subsubsection{チェビシェフ多項式への連続}

$x=\cos\theta$ と書けば,
\[
  T_n(x)=\cos(n\theta)
\]
は
\[
  T_{n+1}(x)=2xT_n(x)-T_{n-1}(x)
\]
を満たす.
三角関数型の漸化式を,\,$\theta$ を消して $x$ の多項式として見直したものが,
チェビシェフ多項式である.

\subsection{パターン9：特性方程式の重解には $n$ を掛ける}

\subsubsection{重解を見たときの予想}

\[
  a_{n+2}-2r a_{n+1}+r^2a_n=0
\]
の特性方程式は
\[
  \lambda^2-2r\lambda+r^2=(\lambda-r)^2=0
\]
であり,根 $r$ が重なっている.
この場合は,
\[
  a_n=(A+Bn)r^n
\]
を予想する.

\subsubsection{なぜ $n$ が現れるのか}

$r\neq0$ として,
\[
  b_n=\frac{a_n}{r^n}
\]
と置く.
元の漸化式を $r^{n}$ で割ると,
\[
  b_{n+2}-2b_{n+1}+b_n=0
\]
となる.
これは二階差分が $0$ の漸化式である.
したがって $b_n$ は等差型,
\[
  b_n=A+Bn
\]
である.
元に戻せば,
\[
  a_n=(A+Bn)r^n.
\]

重解の公式を暗記する代わりに,
\[
  \text{共通の指数 }r^n\text{ を外す}
  \longrightarrow
  \text{差分が0}
\]
という流れで理解できる.

\subsection{パターン10：畳み込みが出たら,母関数を掛ける}

\subsubsection{和の中に $a_ka_{n-k}$ がある場合}

\[
  a_{n+1}=\sum_{k=0}^{n}a_ka_{n-k}
\]
のように,添字が $k$ と $n-k$ に分かれている和を畳み込みという.
この形は,各項を直接追うより母関数を使う方が自然である.

\[
  A(x)=\sum_{n=0}^{\infty}a_nx^n
\]
と置くと,
\[
  A(x)^2
  =\sum_{n=0}^{\infty}
  \left(\sum_{k=0}^{n}a_ka_{n-k}\right)x^n.
\]
畳み込みが積になった.

\subsubsection{例題：カタラン型の漸化式}

\[
  C_0=1,\qquad
  C_{n+1}=\sum_{k=0}^{n}C_kC_{n-k}
\]
を考える.
母関数
\[
  C(x)=\sum_{n=0}^{\infty}C_nx^n
\]
を用いる.
左辺は
\[
  \sum_{n=0}^{\infty}C_{n+1}x^{n}
  =\frac{C(x)-1}{x}
\]
であり,右辺は $C(x)^2$ だから,
\[
  \frac{C(x)-1}{x}=C(x)^2.
\]
すなわち,
\[
  xC(x)^2-C(x)+1=0.
\]
二次方程式として解くと,
\[
  C(x)=\frac{1\pm\sqrt{1-4x}}{2x}.
\]
$C(0)=1$ となる枝を選ぶため,
\[
  C(x)=\frac{1-\sqrt{1-4x}}{2x}.
\]

ここでは,非線形漸化式を項ごとに解かなかった.
無限個の未知数を一つの関数にまとめることで,
漸化式を代数方程式に翻訳したのである.

\subsection{パターン11：一般項を求めず,固定点と不等式を調べる}

\subsubsection{非線形漸化式の見抜き方}

\[
  a_{n+1}=f(a_n)
\]
で $f$ が二次式や分数式になると,
特性方程式のような一般的な公式が使えないことが多い.
そのときは,次の順に調べる.

\begin{enumerate}
  \item 固定点 $L=f(L)$ を求める.
  \item 初期値から動く範囲を確認する.
  \item $f(x)-x$ の符号で増減を調べる.
  \item $|f'(L)|$ で固定点の安定性を予想する.
  \item 単調性と有界性から収束を証明する.
\end{enumerate}

\subsubsection{例題：平方根へ近づく数列}

\[
  a_1>0,\qquad
  a_{n+1}=\frac12\left(a_n+\frac{2}{a_n}\right)
\]
を考える.
固定点 $L$ は
\[
  L=\frac12\left(L+\frac2L\right)
\]
を満たすから,
\[
  L^2=2
\]
であり,正の固定点は $\sqrt2$ である.

さらに,
\[
  \begin{aligned}
    a_{n+1}-\sqrt2
     & =\frac12\left(a_n+\frac2{a_n}-2\sqrt2\right) \\
     & =\frac{(a_n-\sqrt2)^2}{2a_n}.
  \end{aligned}
\]
右辺は常に $0$ 以上である.
また,\,$a_n\ge\sqrt2$ なら
\[
  a_{n+1}\le a_n
\]
も示せる.
よって,\,$\sqrt2$ 以上で単調減少する数列として収束し,
極限は固定点方程式から $\sqrt2$ である.

一般項を得ていなくても,
\[
  \text{どこへ近づくか}
\]
を完全に説明できる.

\subsection{パターン12：型を見抜くための最終チェック}

\subsubsection{式を見たら,次を確認する}

\begin{center}
  \begin{tabular}{c|c}
    式の特徴                & 最初に試す操作             \\ \hline
    $a_n$ に定数倍と定数の和     & 不動点を引く              \\
    $f(n)$ が加わる         & 補助数列を加えて同次化         \\
    $a_{n+1}-a_n$ が既知   & 和を取る                \\
    $a_{n+1}/a_n$ が既知   & 積を取る・望ましい積で割る       \\
    $a_n$ が指数にある        & 対数を取る               \\
    分母に $a_n$ がある       & 逆数・固定点からの比          \\
    $2\cos\theta$ と二階の差 & $\cos(n\theta)$ を試す \\
    特性根が重なる             & $n r^n$ を試す         \\
    $a_ka_{n-k}$ の和     & 母関数を作る              \\
    非線形反復               & 固定点・単調性・有界性         \\
  \end{tabular}
\end{center}

\subsubsection{解法を書くときの最低限の型}

一般項を予想したときは,次の三段階を明示するとよい.

\begin{enumerate}
  \item 初期値で候補が一致すること.
  \item 候補が漸化式を満たすこと.
  \item したがって帰納法,または漸化式の一意性により一致すること.
\end{enumerate}

特に三角関数型では,倍角公式を示して終わりにせず,
「同じ漸化式と同じ初期値を持つ」という最後の一歩を書く.
これによって,思いつきが証明へ変わる.

\subsubsection{解法から理論へ}

ここで紹介した操作は,個別のテクニックとして終わらない.

\[
  \begin{array}{c}
    \text{不動点を引く}          \\
    \downarrow             \\
    \text{固定点と安定性}         \\[2mm]
    \text{補助数列で同次化}        \\
    \downarrow             \\
    \text{線形性・行列・固有値}      \\[2mm]
    \text{三角関数型を見抜く}       \\
    \downarrow             \\
    \text{チェビシェフ多項式・直交多項式} \\[2mm]
    \text{畳み込みを母関数にする}     \\
    \downarrow             \\
    \text{形式的冪級数・組合せ論}     \\[2mm]
    \text{固定点へ近づく数列を作る}    \\
    \downarrow             \\
    \text{ニュートン法・漸近解析}
  \end{array}
\]

というように,問題を解くための置き方が,
より一般的な数学の理論へそのまま伸びていく.
具体的な型を見抜くことは,第4章の抽象理論への入口である.
